\documentclass[11pt]{article}
\usepackage[margin=1in]{geometry}
\usepackage{amsmath,amssymb,amsthm,booktabs}
\usepackage[hidelinks]{hyperref}
\newtheorem{proposition}{Proposition}
\newcommand{\R}{\mathbb{R}}
\title{Relational Verification of Weighted Mixture-of-Experts Programs:\\Scoped Reuse and Explicit Proof Contracts}
\author{Author details omitted for pre-submission review}
\date{Technical review draft, 21 September 2026}
\begin{document}
\maketitle
\begin{abstract}
Input perturbations can change a mixture-of-experts (MoE) program's dispatch
without violating its output property. Requiring route invariance can therefore
exclude safe regions, while independently bounding experts loses common-input
relationships. We develop an ACT analysis using HybridZ for output-layer weighted
MoE programs that retains legal route guards and shared input factors. It
combines expert-wise sufficient conditions, property-directed weighted
obligations, scoped fact reuse, and route-complexity scheduling. In a frozen
comparison on 100 new verification inputs and three fixed same-family models,
the implementation gains 23 policy-accepted safety results over matched
monolithic solving with no losses; all gains involve multiple recorded legal
routes. These are empirical numerical-policy outcomes: a subsequent exact
input audit finds containment gaps, and the historical bounds lack a complete
source-to-output proof. A separate ablation supports the contribution of shared
expert relationships. An ACT-fronted plain CROWN path is much cheaper and has
more numerical positive filters overall, with complementary ACT-only outcomes.
Independent rational checking establishes complete positive results for some
supplied-HZ requests under upstream assumptions. A newer declared-source
construction checks a convolutional enclosure but does not obtain a complete
positive output on its new matrices. We thus contribute a scoped organization
of weighted verification obligations and evidence of its benefits and limits,
not universal backend superiority or source-complete high-accuracy certification.
\end{abstract}

\paragraph{Review status.}
This is a condensed, venue-neutral research draft, not a claim of completed
human review, anonymous artifact release, or acceptance readiness. Detailed
derivations and preserved negative experiments accompany the longer sections.
No new verification experiment was run to produce this manuscript.

\section{Problem and contribution}
Data-dependent routing makes a MoE a piecewise program. The verification
problem concerns all permitted executions, not merely the route observed at
the clean input. An expert output possible somewhere in an input box may be
impossible where that expert is selected. Two individually possible expert
outputs may also be impossible at the same input. Erasing these relationships
can lose useful safety information even after every path has been enumerated.

Our software-analysis question is: \emph{how should a verifier retain dispatch
conditions and common-input relationships while organizing complete weighted
output obligations within a fixed budget?} We study eval-mode, output-layer
selected-softmax top-2 programs on CPU/float64. The normalized top-$k$ algebra
generalizes, but the production interface evaluated here does not cover
arbitrary internal or token-level routing histories.

We contribute (1) a relational organization of legal route and output
obligations using ACT and HybridZ; (2) property-directed weighted verification with
domain-scoped fact reuse and a frozen structural schedule; and (3) an evaluation
that separates empirical policy acceptance, conditional rational proofs, and
checked source construction. Hybrid zonotopes, convexity, McCormick envelopes,
proof reuse, and algorithm selection are not individually claimed as new.

The main comparison shows new-input gains over two internal configurations,
not a population-wide certified accuracy or independent-tool dominance. We
retain the external path's lower cost, the second family's absence of positive
production results, and the source-containment obstruction in the main text.
These findings limit the practical and formal claims of the method.

\section{Semantics and conditional proof rule}
Let $X$ be the \emph{declared} perturbation set, $r:X\to\R^E$ the router,
and $E_i:X\to\R^C$ the experts. A legal unordered route $S$ of size $k$
satisfies $r_i(x)\ge r_j(x)$ for every $i\in S,j\notin S$. The weak
inequalities include every tie-legal route. On its guarded domain $X_S$,
the output is
\[
 F_S(x)=\sum_{i\in S}\lambda_i(x)E_i(x),\qquad
 \lambda_i(x)\ge0,\quad\sum_{i\in S}\lambda_i(x)=1.
\]
For selected-softmax, normalization is over the selected experts, not all
experts. Fix rows $p_\ell(v)=q_\ell^\top v+c_\ell$; classification uses
the clean label's margin against every other class. The target is
\begin{equation}\label{eq:target}
 \forall x\in X,\ \forall S\text{ legal at }x,\ \forall\ell:
 p_\ell(F_S(x))>0.
\end{equation}
This is a mathematical real-program target. The production float64
implementation does not thereby prove its own deployed arithmetic semantics.

\paragraph{Shared relation.}
A hybrid zonotope represents an affine image of bounded continuous and
discrete factors subject to linear constraints. Router and expert analyses
share the input-factor identities and relevant route constraints. Distinct
expert activation choices remain distinct factors. Duplicating input factors
loses correlation; merging independent activation factors can unsoundly
remove executions. Completed route analysis supplies a covering family
$\widehat{\mathcal S}$; the implementation requires its exact-query protocol
to complete under the numerical policy. An incomplete query is not evidence
that the unexamined route is unreachable.

\paragraph{Scoped facts.}
Write $X_i$ for the domain where expert $i$ belongs to some legal route.
Since $X_S\subseteq X_i$ for $i\in S$, a valid membership-domain bound
$p_\ell(E_i(x))\ge L_{i\ell}>0$ remains applicable on $X_S$.
If every selected expert has such a fact for the same row, normalized
nonnegative weights prove the row with lower bound
$\min_{i\in S}L_{i\ell}$. This rule transfers a fact, not an abstract
state with possibly incompatible factor numbering. Identity includes model,
input domain, property, expert, factor frame, and numerical contract.

\paragraph{Residual weighted obligations.}
For pair $S=\{a,b\}$, define in the \emph{same} factor space
\[
 u=q^\top E_b+c,\quad d=q^\top(E_a-E_b),\qquad
 p(F_S)=u+\lambda d.
\]
F0 introduces one scalar product for this property instead of a product for
every output coordinate. Given sound intervals $\lambda\in[l,h]$ and
$d\in[L,H]$, its outer envelope includes
\begin{align*}
 w&\ge l d+L\lambda-lL,& w&\ge h d+H\lambda-hH,\\
 w&\le h d+L\lambda-hL,& w&\le l d+H\lambda-lH.
\end{align*}
The common-input dependence between $u$ and $d$ is retained. The envelope
may still discard dependence of the gate on those factors. It is not an exact
MILP encoding of softmax. Nonpositive relaxed bounds alone establish neither
network unsafety nor impossibility of a stronger proof.

\begin{proposition}[Conditional complete-output composition]\label{prop:complete}
For a top-2 program, suppose $\widehat{\mathcal S}$ covers every tie-legal route on $X$.
For every route/property obligation, suppose either valid scoped expert facts
prove the row, or an outer enclosure of its joint guarded relation has a
valid strictly positive lower bound for the weighted property. If these
facts and enclosures bind the same model, domain, property and semantics,
then~\eqref{eq:target} holds.
\end{proposition}
\begin{proof}
Fix an input, a legal route and a property. Coverage includes that route.
In the reuse case, domain containment makes every expert fact applicable
at this input; convexity gives the minimum positive expert bound, including
the property constant because the weights sum to one. In the residual case,
the concrete shared factors, true gate and $w=\lambda d$ satisfy the outer
enclosure, so its valid lower bound also bounds the concrete property.
Every obligation is discharged and the choice of execution was arbitrary.
\end{proof}
The proposition does not assert that every floating implementation satisfies
its hypotheses. A missing route, unchecked row or expired checker cannot
discharge an obligation. An expert-only violation is UNKNOWN unless an
in-domain input violates the complete routed model on replay.

\section{Execution, reuse and fairness}
The current execution first obtains route coverage and cheap guarded expert
interval facts, without support-tightening solves in that common prelude.
With one legal pair it goes directly to weighted solving: gate weights remain
input-dependent. With multiple pairs, it allocates at most 25\% of the then
remaining budget to expert-wise verification before residual F0. Incomplete
Tier~1 solving may continue to F0 while total budget remains. A replayed
full-model violation or a complete proof can terminate early.

The 25\% rule is a frozen design, not an optimal-allocation theorem. Prior
development found monolithic solving superior in total coverage and staged
solving complementary on multi-route cases. That finding motivated the rule;
it is preserved rather than relabeled by later results.

The primary matched-monolithic comparator computes the same common facts
independently and pays their cost. For each property it partitions pairs into
reused and residual obligations, solves the residual disjunction, and takes
the minimum across both parts. The legacy monolithic configuration remains a
secondary comparator; it is not retroactively forced to adopt the new prelude.
These are internal configurations with shared HZ/F0 machinery, not three
independent external tools. Method-specific Tier~1 search is not silently
donated to the other arm as a free fact.

Every request has a 300-second outer cap including startup, loading, route
analysis, propagation, construction and solving. Common-fact snapshots are
published before method-specific search. Terminal ledgers retain watchdog
terminations even when no final package exists. Separate post-run audits
are not charged as if they were production verification; independently
checked evidence experiments separately disclose their own accounting.

\section{Guarantee contracts and independent evidence}\label{sec:contracts}
We distinguish four layers, which must not be collapsed into a single SAFE
claim. First, Proposition~\ref{prop:complete} is a real-arithmetic result
conditional on valid coverage and enclosures. Second, historical HZ/HiGHS
SAFE denotes acceptance by frozen optimal-status, bound-correction and
positive-margin policies. Third, rational checking validates supplied
constraints and lower-bound certificates. Fourth, a declared-source checker
can establish how particular matrices enclose a particular graph and domain.
None automatically proves execution equivalence with arbitrary native kernels.

\paragraph{Concrete obstruction in the main cohort.}
A saved-only exact audit reconstructs the frozen input formula rather than
recovering old intermediate traces. The requested rational $2/255$ box is
not contained in that formula's box on all 100 inputs (257,850 inward
coordinates); interpreting epsilon as its exact binary64 value also fails
on all inputs (197,704 coordinates). Both interpretations have 50,512 outward
coordinates across all 100 inputs: the boxes are not nested. All 23 primary
gain model--input pairs are affected. This is not a network counterexample,
but it prevents a source-complete real-box interpretation of the main table.
Common defects across arms, tiny gaps, or large old output margins do not
repair the missing containment argument. Corrected construction could change
both route coverage and outcomes; no such counterfactual result is reported.

Independently, the native solver correction policy lacks a checked general
error-coverage argument. Its $10^{-9}$ correction scales and $10^{-7}$
candidate-policy tolerances concern different quantities; comparing those
numbers is not a soundness proof or a proof of unsoundness. Downstream
conversion remains another unclosed premise for the historical bounds.

\paragraph{Exact downstream checking.}
For a finite-box LP $\min c^\top z+d$ with $Az\le b$, $Ez=h$ and
$l\le z\le u$, take $y\le0$, arbitrary $v$, and
$r=c-A^\top y-E^\top v$. Then
\[
 B=d+b^\top y+h^\top v+\sum_j\min(r_jl_j,r_ju_j)
\]
is a valid lower bound. This follows by substitution, the signs of $y$,
and coordinatewise minimization of the residual on the box. The checker
uses exact rational arithmetic; a numerical solver only proposes multipliers.
It reconstructs projections and McCormick rows from the supplied shared HZ,
with explicit continuous relaxation of binary factors. It need not trust a
solver's optimal-status or objective report. Validity of the supplied HZ,
guard lowering and route exclusions nevertheless remains an upstream premise.

The older input98 bundle checks all nine output obligations positively under
those premises. Its portable check does not load a checkpoint or solver.
A later path checks represented-input containment, affine compensation,
expert layers, factor maps, route exclusions and new output construction for
a declared convolutional graph. Its new output bounds are nonpositive or
missing. Positive duals for the old matrices cannot be attached to these new
matrices. The complete source-checked positive request remains unestablished.
Input98 follow-up is stopped; this draft initiates no new local optimization.

\section{Evaluation}
We ask four separate questions: RQ1, does the frozen execution improve
complete policy outcomes on new inputs? RQ2, does common-input relationship
retention contribute? RQ3, what value remains against a runnable external
path? RQ4, which guarantees and model families transfer? Different cohorts
and evidence grades are not pooled.

% Generated by render_submission_tables.py; accounting, not reproof.
\begin{table}[ht]
\centering\small
\caption{New-input confirmation. Each cell is policy-accepted SAFE / solved, out of 100. Solved includes replayed UNSAFE. All arms have the same 300-second request cap.}
\label{tab:confirmation}
\begin{tabular}{lrrr}
\toprule
Model & Adaptive & Matched & Legacy \\
\midrule
Seed 0 & 59 / 89 & 50 / 76 & 46 / 68 \\
Seed 1 & 57 / 81 & 47 / 65 & 45 / 63 \\
Seed 2 & 63 / 86 & 59 / 78 & 50 / 68 \\
\bottomrule
\end{tabular}
\end{table}

\begin{table}[ht]
\centering\small
\caption{External-path and cross-family results. P is a HZ-policy acceptance (HZ) or CROWN numerical filter (CF), not a common strict-certificate grade. U is replayed UNSAFE; ? is UNKNOWN; T is TIMEOUT. Costs include unsuccessful requests.}
\label{tab:external}
\begin{tabular}{llrrrrrr}
\toprule
Cohort / method & Grade & N & P & U & ? & T & Mean s \\
\midrule
MLP / ACT & HZ & 30 & 11 & 9 & 5 & 5 & 138.11 \\
MLP / CROWN path & CF & 30 & 13 & 0 & 17 & 0 & 4.23 \\
Conv / Adaptive & HZ & 30 & 0 & 17 & 0 & 13 & 188.31 \\
Conv / Matched & HZ & 30 & 0 & 11 & 0 & 19 & 217.50 \\
Conv / CROWN path & CF & 30 & 1 & 7 & 22 & 0 & 4.40 \\
\bottomrule
\end{tabular}
\end{table}


\paragraph{RQ1: new-input internal comparison.}
The primary family consists of three frozen eight-expert MLP training runs,
approximately 48\% clean accuracy. The 100 images are common clean-correct
inputs selected by dataset order after excluding prior verification sets,
not by route complexity or expected outcome. Earlier whole-test telemetry
is disclosed: these are new verification endpoints, not wholly unseen images.
Each request specifies a clipped $L_\infty$ radius $2/255$;
Section~\ref{sec:contracts} distinguishes this requested set from the
constructed input set.
The earlier 30-input study is separate. Model, configuration, radius and
selection were frozen before the 900 calls; order rotates within input blocks.

Table~\ref{tab:confirmation} totals 179, 156 and 141 policy acceptances.
Relative to matched, adaptive gains 23 SAFE and 37 solved model--input pairs
without losses. Every SAFE gain has multiple recorded policy-feasible pairs;
two finish in Tier~1 and 21 in F0 with reuse recorded. This is not a concrete
route-flip witness for every native implementation, or casewise causal
attribution to reuse. Relative to legacy there are 40 SAFE gains and two
losses, so no solution-set dominance. All 900 calls are retained, with 739
complete packages and 161 outer timeouts; 300 common-fact pairs agree.

The mean primary SAFE difference is 7.67 percentage points with the frozen
input-clustered descriptive bootstrap interval [4.67,11.00]. Resampling keeps
the three fixed models together for each input; 300 pairs are not 300
independent images or a population of architectures. Mean all-request costs
are 81.71, 111.57 and 153.27 seconds. Different result sets and concentrated
long-tail savings preclude a uniform same-proof speedup claim.

\paragraph{RQ2: shared-input ablation.}
On ten observed inputs and three models, a separate 60-call study replaces
the shared expert relation with a block-diagonal product that retains each
expert's constraints but duplicates their input factors. Every original
joint assignment embeds in that product; it is a sound enlargement given
the source abstractions. The shared arm gains three policy SAFE and five
solved pairs without losses. For seed1/index4018 and seed2/index4014, the
independent arm completes but remains relaxation-undecided; both are
multi-route. For seed2/index4018 it reaches a solver limit. This supports a
relational precision mechanism on two cases, not an attribution of all
timing or RQ1 gains. Variable count and solver cost also change.

\paragraph{RQ3: runnable external path.}
The comparator is explicitly ACT route analysis plus plain CROWN on each
variable-weight static pair over the whole input box. Router and experts
still share the same input; weights are not frozen and the clean route is
not the only obligation. Proving every pair on the whole box is a sufficient
path, though stricter than proving it only on its guard. The configuration
is not full alpha-beta-CROWN with branch-and-bound.

On ten observed inputs and three models (Table~\ref{tab:external}), the
external path has 13 numerical positive filters versus 11 HZ-policy
acceptances, with eight shared, three ACT-only and five CROWN-only outcomes.
All ACT-only counterparts complete with nonpositive CROWN bounds, not frontend
errors; two are multi-route. All CROWN-only counterparts are ACT solver-limit
UNKNOWN. Means are 4.23 versus 138.11 seconds, including cross-environment
handoff. Immutable raw-input preparation is disclosed separately and excluded
equally. This establishes limited complementarity and a real cost challenge,
not ACT dominance; grades differ, and counterexample counts are not extra
safety certificates or a matched attack comparison.

\paragraph{RQ4: transfer and proof completion.}
A separately fixed four-expert convolutional family has 155,052 parameters
and 67.06\% full-test accuracy. Checkpoint selection uses validation accuracy,
not verification results. Its 30-input comparison has no HZ-policy positive;
adaptive's six additional decisions are replayed violations. In a separate
20-new-input evidence study, no arm produces a positive result: matched
returns ten UNSAFE and ten TIMEOUT, evidence mode twenty TIMEOUT, and CROWN
seven UNSAFE and thirteen UNKNOWN. A postselected conditional proof therefore
does not establish transferable positive coverage.

The fixed MLP index3000 request has 18/18 conditionally checked output
obligations. Transferring fresh evidence generation to the three ACT-only
external cases yields 46/54 positive obligations but zero complete positive
requests; these cases are postselected, not a new success-rate estimate.
The old input98 conditional proof and the newer source-checked output
noncompletion in
Section~\ref{sec:contracts} further separate checking coverage from positive
coverage. A checked nonpositive lower bound is neither an unsafe witness
nor proof that the LP optimum cannot be positive.

\section{Related work and novelty boundary}
Hybrid zonotopes already represent ReLU networks~\cite{hz}; our contribution
is not that representation alone. DNNV supports property reduction and
cross-verifier interfaces~\cite{dnnv}. We specialize obligation formation to
tie-legal dispatch, normalized weighted outputs and shared-factor identity.
MetaMoE studies compositional MoE verification~\cite{metamoe}; our
route-invariance comparator is a labeled reimplementation, not execution of
its authors' verifier. Route changes not necessarily implying output errors
are not claimed as a new observation.

CROWN provides static-network bounds~\cite{crown}; modern constraint-driven
methods such as Clip-and-Verify exploit constraints within branch-and-bound
~\cite{clip}. Our plain-CROWN experiment neither measures those complete
systems nor sets their capability ceiling. Our monolithic F0 is a disjunction
of the same outer relaxations, not exact softmax MILP or a public-tool ranking.

FastCert reuses intermediate templates across queries~\cite{fastcert};
incremental Marabou verification inherits conflicts under query refinement
~\cite{conflicts}. We instead reuse row-specific expert facts within a single
request using pair-to-membership domain containment. We do not claim first
use of refinement or proof sharing in neural verification. Source-semantics
frameworks such as TorchLean~\cite{torchlean} address a related semantic gap;
our executed rational checker is not a Lean-kernel proof of the entire stack.
No speedup from a different task is imported as a baseline result here.

\section{Limitations and reproducibility}
The strongest empirical gain is within three fixed, low-accuracy same-family
models on selected clean-correct inputs. The second-family negative results
remain part of the claim, and high-accuracy cross-family route-changing strict
certification is open. Program semantics cover output-layer routing, not
arbitrary token dispatch or co-batch-dependent training-mode BatchNorm.

The exact input audit is a concrete obstruction, not merely a generic
floating-point caveat. Neither structural audits nor source-code labels
establish the hypotheses of Proposition~\ref{prop:complete} for the main
table. Conversely, this does not erase the observed execution differences or
prove the underlying networks unsafe. Upgrading guarantees requires a newly
identified same-source construction and evidence, not relabeling old counts.

Artifacts separate archived-table reconstruction, actual model verification,
and stored-proof checking. Committed review files suffice to rebuild the
tables without numerical packages. A local portable real conditional bundle
has been relocated and checked, but trained weights, raw inputs and full
empirical bundles are not distributed with that table-only workflow. A fresh
source-defined demonstration is a control, not empirical reproduction.
Clean-environment trained-model reproduction, access/licensing arrangements
and human independent method review remain open. No anonymous/public release
or completed external review is asserted by this draft.

\paragraph{AI assistance and author responsibility (disclosure draft).}
AI assistants participated in implementation, tests, experiment orchestration,
analysis scripts, documentation and manuscript preparation. Another AI session
provided criticism; it is not an independent human mathematical review.
Claims in this paper must be traced to recorded executions or stated proofs,
not assistant assertions. Before submission the human authors must finalize
the tool/version inventory, verify the scope of assistance and references,
and take responsibility for the claims. This paragraph is not a statement
that those final human checks have already occurred.

\section{Conclusion}
Retaining dispatch conditions and common-input expert relationships gives a
useful organization of weighted MoE verification obligations. Scoped reuse
and structural scheduling improve frozen-policy outcomes against registered
internal alternatives on new inputs, and a direct ablation identifies a
relational precision benefit. A cheaper external path, negative convolutional
transfer, and an exact source-containment obstruction limit stronger claims.
Conditional positive proofs and checked source construction are both real
progress, but have not yet joined into the claimed stronger positive guarantee.
The result is a scoped verification method and an explicit evidence contract,
not a completed source-to-deployment certification solution.

\begin{thebibliography}{9}
\bibitem{hz} J. Ortiz, A. Vellucci, J. Koeln, and J. Ruths.
Hybrid Zonotopes Exactly Represent ReLU Neural Networks. 2023.
\url{https://arxiv.org/abs/2304.02755}.
\bibitem{dnnv} D. Shriver, S. Elbaum, and M. B. Dwyer.
DNNV: A Framework for Deep Neural Network Verification. 2021.
\url{https://arxiv.org/abs/2105.12841}.
\bibitem{metamoe} Q. Pham, B. Wooding, L. Nam, S. Sasaki, and T. T. Johnson.
MetaMoE: Formal Verification of Compositional Robustness and Scalability of
Mixture-of-Experts Architecture. SAIV 2026 proceedings, pp. 167--190.
\url{https://doi.org/10.1007/978-3-032-32357-6_8}.
\bibitem{crown} H. Zhang, T.-W. Weng, P.-Y. Chen, C.-J. Hsieh, and L. Daniel.
Efficient Neural Network Robustness Certification with General Activation
Functions. NeurIPS 2018. \url{https://arxiv.org/abs/1811.00866}.
\bibitem{clip} D. Zhou, J. Chavez, H. Chen, G. A. Hanasusanto, and H. Zhang.
Clip-and-Verify: Linear Constraint-Driven Domain Clipping for Accelerating
Neural Network Verification. NeurIPS 2025.
\url{https://doi.org/10.52202/085713-5815}.
\bibitem{fastcert} K. Das, S. Ugare, B.-Y. E. Chang, S. Misailovic,
G. Singh, and M. Sridharan.
Uncovering the Limits of Proof Sharing for Neural Networks. 2026.
\url{https://arxiv.org/abs/2608.19351}.
\bibitem{conflicts} R. Elsaleh, L. Davis, H. Wu, and G. Katz.
Incremental Neural Network Verification via Learned Conflicts. 2026.
\url{https://arxiv.org/abs/2603.12232}.
\bibitem{torchlean} R. J. George et al.
TorchLean: Formalizing Neural Networks in Lean. 2026.
\url{https://arxiv.org/abs/2602.22631}.
\end{thebibliography}
\end{document}
